% The two GPIO registers a pin toggle can go through. Atomicity first, speed
% second. Bit 31 is leftmost, as in the reference manual.
\begin{tikzpicture}[font=\sffamily\scriptsize]

% ---------------- the set and reset register
\begin{scope}[yshift=26mm]
\node[regname] at (-1mm,3mm) {GPIOx\_BSRR};
\bitscale
\bitfield{31}{16}{regset}{BR15 to BR0, write 1 to clear that pin}
\bitfield{15}{0}{regfield}{BS15 to BS0, write 1 to set that pin}
\node[regreset, anchor=north west] at (0mm,-0.5mm)
  {write only, reset value 0x0000 0000, one store changes one pin};
\end{scope}

% ---------------- the output data register
\begin{scope}[yshift=6mm]
\node[regname] at (-1mm,3mm) {GPIOx\_ODR};
\bitscale
\bitfield{31}{16}{regrsvd}{reserved}
\bitfield{15}{0}{regfield}{OD15 to OD0, the current output level}
\node[regreset, anchor=north west] at (0mm,-0.5mm)
  {read and write, so a toggle here is read, modify, write: three accesses};
\end{scope}

% ---------------- what the difference actually is
\node[umlnote, text width=60mm, anchor=north west] at (0mm,-6mm)
  {One store to the set and reset register changes a pin with no read at all.
   Nothing else in the port is touched and there is no window between reading
   and writing.};
\node[umlnote, text width=60mm, anchor=north west] at (68mm,-6mm)
  {A read, modify and write of the output register produces the same waveform
   and is not atomic. An interrupt landing inside the window writes back a stale
   value for every other pin in the port.};

\node[font=\sffamily\tiny, text width=128mm, anchor=north west] at (0mm,-26mm)
  {So the difference between the two rungs of the ladder is a correctness
   difference before it is a speed difference, and which one you get depends on
   which call you make. That is the strongest argument available for knowing
   what a layer does, and it is a better argument than any cycle count.};
\end{tikzpicture}
